\section{Introduction}
\label{sec:intro}
Chest CT resolves in three dimensions the anatomy that a radiograph superimposes on one plane, and it can therefore supervise radiograph
models in ways that radiographs alone cannot. The bone and soft-tissue components of a projected CT are known exactly, and findings can be
learned from CT without the overlap of ribs, heart and mediastinum. Both uses are now practical with public data, through projections
rendered from chest CT \citep{angaitkar2026anatomy} and through CT volumes paired with report-derived findings \citep{hamamci2026ctrate}.
Yet CT-supervised models of radiographs are judged mostly by image fidelity or classification accuracy, and these measures hide three
clinically relevant failures: a model may erase a lesion, report a finding that is absent, or justify its output with a map that points
elsewhere. This paper builds CT-supervised models of chest radiographs and tests them for all three.

Ribs and clavicles overlap the lung fields and can obscure abnormalities. Dual-energy subtraction (DES) separates bone from soft tissue but
needs dedicated hardware and a second exposure \citep{kuhlman2006des,sun2026bsldm}. Software bone suppression avoids both; radiologists find
more nodules with suppressed images \citep{freedman2011lung}, and suppressed images have been compared with computer-aided detection
\citep{schalekamp2014cad} and used as classifier inputs \citep{rajaraman2021bone}. Current learning-based suppressors either render paired
supervision from chest CT, projecting bone, non-lung soft tissue and lung separately \citep{angaitkar2026anatomy}, or train diffusion models
on real DES pairs \citep{chen2024bsdiff,sun2026bsldm,sun2025gllcm,sun2026debonedit}. In the CT-supervised line, lesion preservation has been
assessed only indirectly, through detection gains and the contrast of boxed lesions \citep{angaitkar2026anatomy}.

Suppressors are judged by how closely their output matches a bone or soft-tissue reference. A network rewarded for predicting bone can,
however, move lesion signal into the bone image, so that the lesion disappears from the suppressed radiograph; and because a model that does
nothing never absorbs a lesion, absorption has to be read next to suppression strength. Simulated nodules have been used to measure contrast
after hand-built clavicle suppression and in phantoms \citep{hogeweg2013suppression,takagi2020dose}, and learned image-to-image models are
known to remove or invent content in other modalities
\citep{cohen2018distribution,antun2020instabilities,bhadra2021hallucinations,tivnan2024hallucination}. In our search (PubMed, arXiv and
Crossref) we found no audit of learned bone suppression against an exactly known, additive lesion signal, and no training objective aimed
at it.

Whether suppression helps a classifier is a separate question. The suppressed images are computed from the radiograph and add no
information that a classifier could not extract itself; a gain must come from the prior the suppressor learned from CT or from the way a
network uses the extra images. Because a network that reads three images spends three times the computation, a gain can be credited to
suppression only against a control of the same capacity, and seed variation and analytic flexibility can easily produce spurious
differences \citep{varoquaux2022failures,simmons2011falsepositive}.

CT can also teach a radiograph classifier directly. Knowledge distillation trains a student to match a teacher's softened outputs
\citep{hinton2015distilling}, and cross-modal distillation transfers supervision between modalities through paired images
\citep{gupta2016crossmodal}. Because each Qure.ai projection is built from one CT-RATE volume, a CT classifier can teach an X-ray model on
images of the same patients, and one transformer family \citep{liu2022swinv2} can serve both modalities, so that their heat maps can be
compared in the same patient. Whether such a transfer improves real-radiograph classification, aligns the two models' explanations or
reduces hallucinated findings has not been tested.

We address these questions in three studies with analysis plans that were hash-frozen before testing. Study~I measures and reduces lesion
erasure by CT-supervised bone suppressors, Study~II tests whether classifiers gain from the suppressed images, and Study~III distils a Swin
CT teacher into Swin radiograph classifiers. Our contributions are:
\begin{enumerate}
\item a counterfactual lesion audit that measures the fraction of an inserted lesion moved into the bone image, at sites shared by all
 models and reported beside suppression strength;
\item a twin lesion-invariance regulariser, with a dose-response over its weight and an ablation of its parts, trained on public
 CT-derived projections in a three-model ensemble with split-conformal intervals;
\item a pre-specified evaluation on locked sets (paired dual-energy images, an annotated rib set, real nodules, downstream detection and
 classification) and an external validation on paediatric radiographs, with every deviation disclosed;
\item a pre-specified test of bone-suppressed classifier inputs (five configurations including a multi-stream network and a
 capacity-matched control, 175 runs on three datasets) with per-finding analyses, attribution maps scored against expert boxes and
 conformal uncertainty;
\item a cross-modality framework that distils a 2.5D Swin CT teacher into Swin X-ray classifiers through paired projections, evaluated for
 accuracy against a convolutional reference, CT--X-ray heat-map agreement, localisation and hallucination under perturbations that add no
 disease.
\end{enumerate}
Fig.~\ref{fig:workflow} summarises the three studies and Fig.~\ref{fig:architecture} the models.

\begin{figure}[htbp]
\centering
\safefig{figures/fig_workflow.pdf}
\caption{Workflow of the three studies, from data to models, evaluation and the main result. All hypotheses were fixed in analysis plans that
were hash-frozen before testing; training ran in cloud GPU kernels, and every reported number is computed by scripts from the saved
outputs.}
\label{fig:workflow}
\end{figure}

Section~\ref{sec:related} reviews related work and Section~\ref{sec:data} the data. Sections~\ref{sec:methods} to~\ref{sec:audit}
describe the suppressors and the audit, Section~\ref{sec:classifier} the classifier study and Section~\ref{sec:experiments} the
experiments of Studies~I and~II; Section~\ref{sec:xm} presents the cross-modality study, Section~\ref{sec:results} reports all results,
and Sections~\ref{sec:discussion} to~\ref{sec:conclusion} discuss them.
